\chapter{Drawing 3D objects}

\section{The shapes that come in the box}

raylib has no 3D modelling tools, but it has enough primitives to build a whole
game out of, and every one of them is a single call:

\begin{lstlisting}
DrawCube(centre, width, height, length, colour);     // sizes are X, Y, Z
DrawCubeV(centre, size, colour);                     // same, size as Vector3
DrawCubeWires(centre, w, h, l, colour);              // just the edges
DrawSphere(centre, radius, colour);
DrawCylinder(bottomCentre, radiusTop, radiusBottom, height, sides, colour);
DrawPlane(centre, size2D, colour);                   // a flat XZ quad
DrawLine3D(from, to, colour);
DrawGrid(slices, spacing);                           // the XZ floor grid
\end{lstlisting}

A few things worth knowing before you use them:

\begin{itemize}[leftmargin=*,itemsep=3pt]
  \item \lstinline|DrawCube| takes the \textbf{centre}, like
        \lstinline|DrawCircle| did in 2D and unlike \lstinline|DrawRectangle|.
  \item \lstinline|DrawCylinder| takes the \textbf{bottom} centre, not the middle.
        Give it a top radius of 0 and you have a cone.
  \item \lstinline|DrawPlane| is always flat on the XZ plane. It takes a
        \lstinline|Vector2| size because a flat thing only has two dimensions.
  \item \lstinline|DrawGrid| is centred on the origin, so a grid over a level
        that starts at 0 and grows positive will only cover a quarter of it.
\end{itemize}

\section{Why everything looks flat}

The first 3D scene anybody draws looks like coloured paper cut-outs. Every face
of the cube is exactly the same shade, so the shape reads as a hexagon rather
than a box.

That is not a bug. Plain \lstinline|DrawCube| does no lighting at all --- there
is no light in the scene, so every surface just gets its flat colour. Real
lighting means shaders, which is a much later chapter of your life.

Three cheap tricks get you 90\% of the way, and the lessons use all three:

\begin{enumerate}[leftmargin=*,itemsep=3pt]
  \item \textbf{Draw the wireframe on top.} A cube plus its dark edges reads as a
        solid object immediately:
\begin{lstlisting}
DrawCube(p, 2, 2, 2, MAROON);
DrawCubeWires(p, 2, 2, 2, (Color){ 30, 26, 24, 255 });
\end{lstlisting}

  \item \textbf{Vary the colour per object.} Identical walls merge into one mass;
        a little variation separates them. The mini Doom shades each wall cell
        from its grid coordinates:
\begin{lstlisting}
int shade = 88 + ((col*7 + row*13) % 26);
\end{lstlisting}

  \item \textbf{Make the floor and ceiling different.} Two
        \lstinline|DrawPlane| calls in different colours do more for the sense of
        being inside a room than any amount of geometry.
\end{enumerate}

\section{A fake shadow}

Nothing tells you where an object is in 3D like a shadow, and a convincing one
costs one line: a very flat dark cube on the floor, directly under the object.

\begin{lstlisting}
DrawCube((Vector3){ obj.x, 0.01f, obj.z }, 1.0f, 0.02f, 1.0f, Fade(BLACK, 0.5f));
\end{lstlisting}

The \lstinline|0.01f| lifts it a hair above the floor. Put it at exactly
\lstinline|0| and it fights with the floor plane for the same pixels, which
produces a horrible flickering called \emph{z-fighting}. Whenever two surfaces
occupy the same plane and shimmer, nudge one of them.

\section{Making something face the player}

Enemies in early first person games were flat images that always turned to face
you. raylib does that with billboards:

\begin{lstlisting}
DrawBillboard(camera, texture, position, size, WHITE);
\end{lstlisting}

If you have no texture, you can still get the effect with geometry. Work out the
direction from the object to the camera, and place details along it --- this is
the trick the mini Doom uses to give its enemies eyes that follow you:

\begin{lstlisting}
Vector3 toPlayer = Vector3Normalize(Vector3Subtract(camera.position, e->position));
Vector3 side = Vector3Normalize(Vector3CrossProduct(toPlayer, (Vector3){ 0, 1, 0 }));
Vector3 base = Vector3Add(e->position, Vector3Scale(toPlayer, 0.22f));
DrawSphere(Vector3Add(base, Vector3Scale(side,  0.10f)), 0.055f, YELLOW);
DrawSphere(Vector3Add(base, Vector3Scale(side, -0.10f)), 0.055f, YELLOW);
\end{lstlisting}

That is the first time we have used the cross product for something practical.
Chapter 12 explains it properly; for now, read it as ``give me the direction
perpendicular to these two''.

\section{Models, when you outgrow cubes}

\begin{lstlisting}
Model model = LoadModel("robot.glb");        // once, outside the loop
DrawModel(model, position, scale, WHITE);    // every frame
DrawModelEx(model, position, rotAxis, rotAngle, scaleV, WHITE);
UnloadModel(model);                          // once, after the loop
\end{lstlisting}

Same load-outside-unload-after rule as textures. raylib reads \file{.obj},
\file{.gltf}, \file{.glb}, \file{.iqm} and \file{.vox}, and \file{.glb} is the
one to prefer because it packs the mesh, the materials and the textures into a
single file.

\section{The maths helpers}

Vector maths lives in a separate header you have to include yourself:

\begin{lstlisting}
#include "raymath.h"
\end{lstlisting}

The handful you will actually use:

\begin{center}
\begin{tabular}{@{}ll@{}}
\toprule
\textbf{Function} & \textbf{What it gives you} \\
\midrule
\lstinline|Vector3Add(a, b)| & a + b \\
\lstinline|Vector3Subtract(a, b)| & a $-$ b: the vector from b to a \\
\lstinline|Vector3Scale(v, f)| & v stretched by f \\
\lstinline|Vector3Length(v)| & how long it is \\
\lstinline|Vector3Normalize(v)| & same direction, length exactly 1 \\
\lstinline|Vector3Distance(a, b)| & distance between two points \\
\lstinline|Vector3CrossProduct(a, b)| & a direction perpendicular to both \\
\bottomrule
\end{tabular}
\end{center}

\keyidea{\lstinline|Vector3Normalize| is the one to understand. A normalized
vector has length 1, so it carries a \emph{direction} and nothing else. Multiply
it by a speed and you get a velocity; multiply it by a distance and you get an
offset. Almost every movement line in the next chapters is
\lstinline|normalize(something) * speed * dt|.}
